\section{Materials and Methods}
\label{sec:materials_and_methods}

Building on the limitations highlighted in prior studies, this work introduces a simple yet effective classification framework for detecting abnormal gait patterns. 
The approach relies on just four carefully selected features derived from gyroscope signals—extremal values that are both biomechanically meaningful and easy to compute. 
Designed with practical deployment in mind, the framework prioritizes efficiency and interpretability, making it well-suited for use in wearable systems. To ensure that 
the model generalizes well across individuals, particularly in diverse clinical populations, we employ a leave-one-out cross-validation (LOOCV) strategy.

\subsection{Dataset Description}
The dataset used in this study was collected from a total of 18 participants, comprising 9 healthy individuals and 9 stroke survivors. 
The inclusion of both able-bodied and neurologically impaired subjects was intended to support the development of a binary classification 
model capable of distinguishing between normal and abnormal gait patterns. 

**At the time of writing, detailed demographic data (e.g., age, gender, affected limb) are being finalized and will be presented in a subsequent demographics table.

All participants underwent gait data collection using wearable inertial measurement units (IMUs). Each subject was equipped with sensors placed on the left and right 
shanks, as well as the left and right thighs as seen on . These positions were selected to capture comprehensive lower-limb kinematic data, with particular emphasis on the 
angular velocity along the z-axis—corresponding to the rotational movement of the legs during gait. The IMUs recorded tri-axial gyroscope data at a consistent 
sampling rate of 100Hz, and sessions were conducted under supervised walking trials on a flat indoor surface to ensure safety and signal consistency.

\begin{figure}[h]
    \centering
    \includegraphics[width=0.8\textwidth]{Images/.png}
    \caption{Placement of IMUs on the lower limbs during data collection.}
    \label{fig:imu_placement}
\end{figure}

\subsection{Feature Extraction}




\subsection{Classification Models}

\subsection{Evaluation Strategy}

\subsection{Dataset Description}
